
\documentclass{article}
\usepackage{colortbl}
\usepackage{makecell}
\usepackage{multirow}
\usepackage{supertabular}

\begin{document}

\newcounter{utterance}

\centering \large Interaction Transcript for game `cladder', experiment `full\_v1.5\_default', episode 5155 with Qwen3.5{-}4B.
\vspace{24pt}

{ \footnotesize  \setcounter{utterance}{1}
\setlength{\tabcolsep}{0pt}
\begin{supertabular}{c@{$\;$}|p{.15\linewidth}@{}p{.15\linewidth}p{.15\linewidth}p{.15\linewidth}p{.15\linewidth}p{.15\linewidth}}
    \# & \multicolumn{2}{c}{Player} && \multicolumn{2}{c}{Game Master} \\
    \hline

    \theutterance \stepcounter{utterance}  
    & & & \multicolumn{4}{p{0.6\linewidth}}{
        \cellcolor[rgb]{0.9,0.9,0.9}{
            \makecell[{{p{\linewidth}}}]{
                \texttt{\tiny{[P1$\langle$GM]}}
                \texttt{You are an expert in causal inference. The following question is not a typical commonsense query, but rather a meticulously designed question created by a professor specializing in causal inference, intended to assess the students' mastery of the course content.} \\
\texttt{Imagine a self{-}contained, hypothetical world with only the following conditions, and without any unmentioned factors or causal relationships: Talent has a direct effect on elite institution admission status. Brown eyes has a direct effect on elite institution admission status. Method 1: We look at how talent correlates with brown eyes case by case according to elite institution admission status. Method 2: We look directly at how talent correlates with brown eyes in general. To understand how talent affects brown eyes, is it more correct to use the Method 1 than Method 2?} \\
\texttt{Start your answer with 'Yes' or 'No', followed by additional reasoning or evidence to support your explanation.} \\
            }
        }
    }
     \\ \\

    \theutterance \stepcounter{utterance}  
    & \multicolumn{4}{p{0.6\linewidth}}{
        \cellcolor[rgb]{0.561,0.737,0.561}{
            \makecell[{{p{\linewidth}}}]{
                \texttt{\tiny{[P1$\rangle$GM]}}
                \texttt{No} \\
\\ 
\texttt{**Reasoning:**} \\
\\ 
\texttt{In this hypothetical world, we are asked to determine the causal effect of **Talent** on **Brown Eyes**. However, the problem statement explicitly defines the causal structure as:} \\
\texttt{1.  Talent \$\textbackslash{}rightarrow\$ Elite Institution Admission} \\
\texttt{2.  Brown Eyes \$\textbackslash{}rightarrow\$ Elite Institution Admission} \\
\\ 
\texttt{There is **no direct causal link** defined between Talent and Brown Eyes. In causal inference, asking "how Talent affects Brown Eyes" when there is no arrow connecting them implies we are looking for an indirect effect (mediated through admission) or a spurious correlation.} \\
\\ 
\texttt{Let's analyze the two methods:} \\
\\ 
\texttt{*   **Method 2 (General Correlation):** This calculates the total correlation between Talent and Brown Eyes. Since both variables independently cause a third variable (Admission), they will likely be correlated due to **selection bias** (or collider stratification bias). If we observe people who got admitted, we might see a correlation between their talent and eye color simply because both factors "selected" them into the group, not because one caused the other. If we look at the general population, the correlation might be zero or weak if the effects on admission are independent. Crucially, this method measures association, not causation, and does not isolate the specific pathway of interest (which doesn't exist).} \\
\\ 
\texttt{*   **Method 1 (Case{-}by{-}case according to admission status):** This method involves stratifying the data by the common effect (Admission Status) and observing the relationship between Talent and Brown Eyes within each stratum (e.g., among admitted students vs. non{-}admitted students).} \\
\texttt{    *   If we look at the **Admitted** group: We are conditioning on a collider. This often induces a spurious correlation (or strengthens an existing one) between Talent and Brown Eyes, making it appear as though they are related when they are not causally linked.} \\
\texttt{    *   If we look at the **Non{-}Admitted** group: Similarly, conditioning on the non{-}outcome of a collider can induce spurious associations.} \\
\\ 
\texttt{**Conclusion:**} \\
\texttt{Neither method correctly identifies the causal effect of Talent on Brown Eyes because **there is no such causal effect** in the defined world. The question asks "To understand how talent affects brown eyes," implying a causal relationship exists. Since the graph is \$Talent \textbackslash{}rightarrow Admission \textbackslash{}leftarrow Brown Eyes\$, Talent does not affect Brown Eyes.} \\
\\ 
\texttt{If the question implies "Which method best estimates the *spurious* correlation induced by the common effect?", Method 1 (stratification) is actually the mechanism that *creates* the bias (collider stratification), whereas Method 2 (marginal correlation) might show no relationship if the effects are independent. However, in standard causal inference, you do not use stratification on a collider to understand a non{-}existent causal link; you use it to *detect* the bias.} \\
\\ 
\texttt{Therefore, it is **not correct} \\
            }
        }
    }
    & & \\ \\

    \theutterance \stepcounter{utterance}  
    & & & \multicolumn{2}{p{0.3\linewidth}}{
        \cellcolor[rgb]{0.9,0.9,0.9}{
            \makecell[{{p{\linewidth}}}]{
                \texttt{\tiny{[GM$|$GM]}}
                \texttt{No} \\
            }
        }
    }
    & & \\ \\

    \theutterance \stepcounter{utterance}  
    & & & \multicolumn{2}{p{0.3\linewidth}}{
        \cellcolor[rgb]{0.9,0.9,0.9}{
            \makecell[{{p{\linewidth}}}]{
                \texttt{\tiny{[GM$|$GM]}}
                \texttt{no} \\
            }
        }
    }
    & & \\ \\

    \theutterance \stepcounter{utterance}  
    & & & \multicolumn{2}{p{0.3\linewidth}}{
        \cellcolor[rgb]{0.9,0.9,0.9}{
            \makecell[{{p{\linewidth}}}]{
                \texttt{\tiny{[GM$|$GM]}}
                \texttt{game\_result = WIN} \\
            }
        }
    }
    & & \\ \\

\end{supertabular}
}

\end{document}
